% !TeX spellcheck = en_US
\documentclass[french]{yLectureNote}

\title{Introduction à la termodynamique}
\subtitle{Thermodynamique}
\author{Paulhenry Saux}
\date{\today}
\yLanguage{Français}

\professor{}

\usepackage{graphicx}%----pour mettre des images
\usepackage[utf8]{inputenc}%---encodage
\usepackage{geometry}%---pour modifier les tailles et mettre a4paper
%\usepackage{awesomebox}%---pour les boites d'exercices, de pbq et de croquis ---d\'esactiv\'e pour les TP de PC
\usepackage{tikz}%---pour deiffner + d\'ependance de chemfig
\usepackage{tkz-tab}
\usepackage{tabularx}%---pour dimensionner automatiquement les tableaux avec variable X
\usepackage{awesomebox}%---Pour les boites info, danger et autres
\usepackage{menukeys}%---Pour deiffner les touches de Calculatrice
\usepackage{fancyhdr}%---pour les en-t\^ete personnalis\'ees
\usepackage{blindtext}%---pour les liens
\usepackage{hyperref}%---pour les liens (\`a mettre en dernier)
\usepackage{caption}%---pour la francisation de la l\'egende table vers Tableau
\usepackage{pifont}
\usepackage{array}%---pour les tableaux
\usepackage{lipsum}
\usepackage{yFlatTable}
\usepackage{multicol}
\newcommand{\Lim}[1]{\lim\limits_{\substack{#1}}\:}
\renewcommand{\vec}{\overrightarrow}
\newcommand{\bz}{\overline{z}}
\DeclareMathOperator{\card}{card}
\renewcommand{\vec}{\overrightarrow}
\newcommand{\N}[0]{\mathbb{N}}
\newcommand{\R}[0]{\mathbb{R}}
\newcommand{\C}[0]{\mathbb{C}}
\newcommand{\dd}[0]{\mathrm{d}}
\newcommand{\er}{\vec{e_{\rho}}}
\newcommand{\ep}{\vec{e_{\varphi}}}
\newcommand{\pip}{\pi^+}
\newcommand{\pim}{\pi^-}
\newcommand{\mv}{\mathcal{V}}
\DeclareMathOperator{\Div}{Div}
\newcommand{\rot}{\vec{\mathrm{rot}}}
\newcommand{\grad}{\vec{\mathrm{grad}}}
\begin{document}
\setcounter{chapter}{2}
\chapter{Principe de conservation de l'énergie}

On étudie un système fermé $\Omega$.

\section{Énoncé et Conservation}

\begin{theorem}[Principe de conservation de Lavoisier]
    L'énergie est une grandeur conservative, elle ne peut \^etre ni créée, ni détruite.
\end{theorem}

% \begin{theorem}[Énergie de l'Univers]
%     L'énergie d'un système isolé est une constante. L'énergie de l'univers est notée $E_0$. 
%     C'est une grandeur extensive, donc $E_{\Omega} + E_{\bar{\Omega}} = E_0$.
% \end{theorem}

\tipsInfo{Principe de localité}{Pour que la variation d'énergie soit instantanée, le transfert doit s'effectuer physiquement à la frontière $\partial \Omega$ sous forme de flux $\varphi_{\partial \Omega}$.}

\begin{theorem}[Bilan d'énergie]
    $$\frac{\dd E_{\Omega}}{\dd t} = \phi_{\partial \Omega}$$
\end{theorem}

\section{Premier principe de la thermodynamique}

\subsection{Nature de l'énergie et des échanges}
L'énergie totale se décompose en :
$$E_{\Omega} = E_{C} + E_p + U$$
\begin{itemize}
    \item $E_c$ : Énergie cinétique macroscopique.
    \item $E_p$ : Énergie potentielle macroscopique.
    \item $U$ : Énergie interne (microscopique).
\end{itemize}

\begin{definition}[Flux d'énergie]
    On appelle flux d'énergie sous forme de chaleur ($\phi_Q$) tous les flux d'énergie qui ne sont pas mécaniques ($\phi_W$), c'est-à-dire non liés au mouvement de la frontière.
\end{definition}



\subsection{Énoncé du premier principe}

\begin{theorem}[Premier principe de la thermodynamique]
    Pour une transformation entre un état $i$ et un état $f$ :
    $$\Delta_{if} (E_c + E_p + U) = W_{i \to f} + Q_{i \to f}$$
\end{theorem}

\criticalInfo{Utilité}{En pratique, pour un système au repos macroscopique, on a $\Delta U = W + Q$. Ce principe est l'outil majeur pour calculer la chaleur $Q$ échangée, souvent difficile à mesurer directement.}

\section{Capacités thermiques et Enthalpie}

\begin{definition}[Capacité thermique]
    Quantité d'énergie thermique nécessaire pour augmenter la température d'un corps de $\dd T$:
    $$\delta Q = C\dd T$$
\end{definition}

\subsection{Capacité à volume constant ($C_v$)}
À volume constant, $\delta W = 0$, donc $\dd U = \delta Q = C_v \dd T$.
\begin{theorem}[Définition de $C_v$]
    $C_v = \left(\frac{\partial U}{\partial T}\right)_V$.
\end{theorem}

\subsection{L'Enthalpie ($H$) et $C_p$}
Pour une transformation monobare ($P_{ext} = Cst$), le travail est $W = -P_{ext}\Delta V$.

\begin{flalign*}
    Q &= \Delta U - W = \Delta U + P_{ext}\Delta V \\
      &= (U_f + P_{ext}V_f) - (U_i + P_{ext}V_i) \\
      &= \Delta(U + PV)
\end{flalign*}

\begin{theorem}[Enthalpie]
    On définit la fonction d'état extensive Enthalpie par : $H = U + PV$.
    Pour une transformation monobare, la chaleur échangée est égale à la variation d'enthalpie : $Q = \Delta H$.
\end{theorem}

\begin{theorem}[Capacité thermique à pression constante]
    $C_p = \left(\frac{\partial H}{\partial T}\right)_P$.
\end{theorem}

\section{Lois de Joule et Gaz Parfaits}

\begin{theorem}[Première loi de Joule]
    Pour un Gaz Parfait (GP), l'énergie interne ne dépend que de la température : $U = U(T)$.
    Par extension, $H = H(T)$ (Seconde loi de Joule).
\end{theorem}

\subsection{Relation de Mayer et Gamma}
\begin{theorem}[Relations fondamentales]
    $$C_p - C_v = nR \quad \text{et} \quad \gamma = \frac{C_p}{C_v}$$
\end{theorem}

\warningInfo{Résolution de système}{En combinant ces deux relations, on obtient les expressions indispensables :
$$C_v = \frac{nR}{\gamma-1} \quad ; \quad C_p = \frac{nR\gamma}{\gamma-1}$$}

\section{Transformations Adiabatiques (Lois de Laplace)}

\criticalInfo{Hypothèses de Laplace}{
    Pour utiliser les relations suivantes, il faut impérativement :
    \begin{itemize}
        \item Un Gaz Parfait (GP).
        \item Une transformation Adiabatique ($Q=0$).
        \item Une transformation Quasi-Statique (QS).
    \end{itemize}
}

\begin{theorem}[Lois de Laplace]
    $$PV^{\gamma} = Cst \quad ; \quad TV^{\gamma-1} = Cst \quad ; \quad P^{1-\gamma}T^{\gamma} = Cst$$
\end{theorem}

\explanation{a1}{Le travail pour une adiabatique QS se calcule par $\Delta U$ ou par intégration de $-P dV$ en remplaçant $P$ par $Cst/V^\gamma$.}

\begin{flalign*}
    W_{i\to f} &= \Delta U = C_v \Delta T = \frac{nR}{\gamma-1}(T_f - T_i) = \frac{P_f V_f - P_i V_i}{\gamma-1} \explain{a1}{right}{0}{0.5}{}
\end{flalign*}



\section{Expériences de Détente}

\subsection{Détente de Joule et Gay-Lussac}
Détente d'un gaz dans le vide (adiabatique et sans travail).
\checkInfo{Résultat}{Comme $W=0$ et $Q=0$, alors $\Delta U = 0$. Pour un GP, la température reste constante ($T_i = T_f$).}

\subsection{Détente de Joule-Thomson}
Écoulement permanent à travers un corps poreux (ou étranglement).
\checkInfo{Résultat}{La transformation est isenthalpique : $\Delta H = 0$. Pour un GP, cela implique également $T_i = T_f$.}
\section{Difficultés}

\subsection{Confusion entre $P$, $P_{ext}$ et $P_{atm}$}
C'est l'erreur la plus fréquente dans le calcul du travail $W = \int -P_{ext} \dd V$.
\begin{itemize}
    \item \textbf{Transformation brutale :} $P_{ext}$ est la pression finale d'équilibre ($P_f$). Elle est constante.
    \item \textbf{Transformation quasi-statique :} On considère $P_{ext} = P_{gaz}$ à chaque instant. On peut alors utiliser la loi des gaz parfaits pour remplacer $P$ dans l'intégrale.
    \item \textbf{Piston à l'équilibre :}  $P_{ext}$ peut inclure le poids du piston : $P_{ext} = P_{atm} + \frac{m_p g}{S_p}$.
\end{itemize}

\subsection{Diagramme de Clapeyron}
Le diagramme $(P, V)$ permet de visualiser le travail comme l'aire sous la courbe.
\begin{itemize}
    \item \textbf{Chemin :} Pour aller d'un point A à un point B, l'aire est plus grande pour une transformation monobare (rectangle) que pour une isotherme (arc d'hyperbole).
    \item \textbf{Cycle :} Si le cycle est parcouru dans le sens horaire, le travail total est négatif (moteur). S'il est anti-horaire, il est positif (récepteur).
\end{itemize}

\subsection{Différencier $\Delta U$ et $\Delta H$}
\tipsInfo{Astuce de choix}{
    \begin{itemize}
        \item Utiliser $\Delta U = W + Q$ pour les transformations à \textbf{volume constant} (isochores).
        \item Utiliser $\Delta H = Q$ pour les transformations \textbf{monobares} (où $W$ est déjà "absorbé" dans la définition de $H$).
    \end{itemize}
}
\end{document}


